% Use the existing owner-derived source mechanism; do not edit frozen originals.
\OLSADeclareDocumentTextCorrection{OLP-0014}{FA0014-HYPHENATION}
{\emph{پاد\-مت\-قارن}}
{\emph{پادمتقارن}}
\OLSADeclareDocumentTextCorrection{OLP-0014}{FA0014-IRREFLEXIVE-HEADING}
{\begin{defn}[غیربازتابی‌بودن]}
{\begin{defn}[پادبازتابی‌بودن]}
\OLSADeclareDocumentTextCorrection{OLP-0014}{FA0014-IRREFLEXIVE-DEFINITION}
{\emph{غیربازتابی}}
{\emph{پادبازتابی}}
\OLSADeclareDocumentTextCorrection{OLP-0014}{FA0014-CONNECTIVITY-EXPLANATION}
{آنگاه دست‌کم یکی از $Rxy$ و~$Ryx$ برقرار است.
\end{defn}}
{آنگاه دست‌کم یکی از $Rxy$ و~$Ryx$ برقرار است.
به بیان دیگر، هر دو عضو متمایز در دامنه با این رابطه مقایسه‌پذیرند.
\end{defn}}
\OLSADeclareDocumentTextCorrection{OLP-0014}{FA0014-IRREFLEXIVE-PARAGRAPH}
{توجه کنید که اگر $A \neq \emptyset$، آنگاه هیچ رابطهٔ غیربازتابی
روی~$A$ بازتابی نیست و هر رابطهٔ نامتقارن روی~$A$ پادمتقارن نیز هست.
بااین‌حال، روابطی $R \subseteq A^2$ وجود دارند که نه بازتابی‌اند و نه
غیربازتابی، و روابط پادمتقارنی نیز وجود دارند که نامتقارن نیستند.}
{توجه کنید که اگر $A \neq \emptyset$، آنگاه هیچ رابطهٔ پادبازتابی
روی~$A$ بازتابی نیست و هر رابطهٔ نامتقارن روی~$A$ پادمتقارن نیز هست.
بااین‌حال، اگر دامنه دست‌کم دو عضو داشته باشد، روابطی $R \subseteq A^2$
وجود دارند که نه بازتابی‌اند و نه پادبازتابی. همچنین، روی هر دامنهٔ
ناتهی رابطه‌ای پادمتقارن وجود دارد که نامتقارن نیست.
\emph{یادداشت ویراستاری:} قید اندازهٔ دامنه برای روشن‌کردن ادعای متن مبدأ
افزوده شده است؛ روی دامنهٔ تک‌عضوی هر رابطه یا بازتابی است یا پادبازتابی.}
\OLSADeclareDocumentTextCorrection{OLP-0015}{FA0015-TWO-ELEMENTS}
{گفته می‌شود !!^{element}s $x$
و $y$}
{گفته می‌شود دو !!{element} $x$
و $y$}
\OLSADeclareDocumentTextCorrection{OLP-0015}{FA0015-MODULAR-SCOPE}
{برای هر
سه عدد صحیح مثبت $a$، $b$ و $n \in \PosInt$}
{برای هر دو عدد طبیعی $a$، $b$ و هر $n \in \PosInt$}
\OLSADeclareDocumentTextCorrection{OLP-0015}{FA0015-COUNTED-NOUN}
{دارای $n$ !!{element}s است}
{دارای $n$ !!{element} است}
\OLSADeclareDocumentTextCorrection{OLP-0015}{FA0015-DISCLOSE-SCOPE}
{یعنی~$\equivrep{n-1}{\equiv_n}$.}
{یعنی~$\equivrep{n-1}{\equiv_n}$.
\emph{یادداشت ویراستاری:} در این مثال، صفر نیز برای دو عدد نخست مجاز دانسته شده است تا دامنه با خارج‌قسمت اعداد طبیعی و نمایندهٔ صفر در ادامه سازگار باشد؛ پیمانه همچنان مثبت است.}
\OLSADeclareSetup{OLP-0014}{\olsa_source_correction_install_document_text:}
\OLSADeclareSetup{OLP-0015}{\olsa_source_correction_install_document_text:}
